\renewcommand{\tema}{Lección: Cinemática}


\section{MAGNITUDES ESCALARES Y VECTORIALES}

\begin{seccion}
    \textbf{Magnitud escalar:} Cantidad que se determina por un número y una unidad.
    \textit{Ejemplos:} masa (5 kg), tiempo (3 s), temperatura (25°C), distancia (10 m)
    \vspace{0.3cm}
    \textbf{Magnitud vectorial:} Cantidad que requiere magnitud, dirección y sentido.
    \textit{Ejemplos:} velocidad (30 m/s hacia el norte), fuerza (50 N hacia arriba), desplazamiento (10 m al este)
    \vspace{0.3cm}
    \textbf{Notación vectorial:}
    $\vec{v} \text{ (con flecha para vectores)}$
    $v \text{ o } |\vec{v}| \text{ (magnitud del vector)}$
\end{seccion}

\subsection {Esquema conceptual}

\begin{center}
\begin{tikzpicture}[scale=1.2]
% Título
\node[draw, fill=azuloscuro, text=white, rectangle, rounded corners, minimum width=8cm, minimum height=0.8cm] at (0,4) {\textbf{MAGNITUDES}};

% Escalar
\node[draw, fill=fondogris, rectangle, rounded corners, minimum width=3.5cm, minimum height=0.7cm] at (-2.5,2.5) {\textbf{ESCALARES}};
\draw[->, thick] (0,3.6) -- (-2.5,3);

\node[align=center] at (-2.5,1.5) {Solo magnitud\\y unidad};
\node[align=center, font=\small] at (-2.5,0.3) {Ejemplos: masa 5 kg, tiempo 3 s};

\node[draw, fill=fondogris, rectangle, rounded corners, minimum width=3.5cm, minimum height=0.7cm] at (2.5,2.5) {\textbf{VECTORIALES}};
\draw[->, thick] (0,3.6) -- (2.5,3);

\node[align=center] at (2.5,1.5) {Magnitud, direccion y sentido};
\node[align=center, font=\small] at (2.5,0.3) {Ejemplos: vel. 30 m/s al norte, fuerza 50 N arriba};
\end{tikzpicture}
\end{center}

\section{DESPLAZAMIENTO VS TRAYECTORIA}
\begin{seccion}

\begin{definicion}
    \textbf{Trayectoria (distancia):} Longitud del camino recorrido. Magnitud escalar.
\end{definicion}

\begin{ecuacion}
    $d = \text{suma de las distancias recorridas}$
\end{ecuacion}

\begin{definicion}
    \textbf{Desplazamiento:} Vector del punto inicial al punto final. Magnitud vectorial.
\end{definicion}

\begin{ecuacion}
$\vec{\Delta x} = \vec{x_f} - \vec{x_0}$

En una dimensión:
$\Delta x = x_f - x_0$
\end{ecuacion}

\textbf{Ejemplo:}
\begin{itemize}
    \item Objeto que da una vuelta: distancia = perímetro, desplazamiento = 0
    \item El desplazamiento puede ser cero, negativo o positivo
    \item La distancia es positiva o cero
\end{itemize}
\end{seccion}

\subsection{Esquema conceptual}

\begin{center}
\begin{tikzpicture}[scale=1.0]
% Camino curvo
\draw[thick, azuloscuro] (0,0) .. controls (1,1) and (2,1) .. (3,0.5) .. controls (4,0) and (4.5,1) .. (5,0);

% Puntos inicial y final
\filldraw[azuloscuro] (0,0) circle (3pt) node[below] {$x_0$};
\filldraw[azuloscuro] (5,0) circle (3pt) node[below] {$x_f$};

% Desplazamiento (flecha recta)
\draw[->, very thick, verdecorrecto] (0,-0.8) -- (5,-0.8) node[midway, below] {Desplazamiento: $\Delta x$};

% Distancia (etiqueta en el camino)
\node[align=center, azuloscuro] at (2.5,1.5) {Distancia = longitud\\del camino};

% Etiquetas
\node[align=center, below] at (2.5,-1.5) {\textbf{Escalar:} $d$ (distancia)\\
\textbf{Vectorial:} $\vec{\Delta x}$ (desplazamiento)};
\end{tikzpicture}
\end{center}


\section{VELOCIDAD VS RAPIDEZ}

\begin{seccion}
\begin{definicion}
    \textbf{Rapidez:} Magnitud escalar que mide qué tan rápido se mueve un objeto.
\end{definicion}

\begin{ecuacion}
\textbf{Ecuación de rapidez:}
$r = \frac{d}{t}$
donde $d$ es la distancia recorrida.
\end{ecuacion}

\begin{definicion}
    \textbf{Velocidad:} Magnitud vectorial que mide el cambio de posición.
\end{definicion}

\begin{ecuacion}
\textbf{Ecuación de velocidad:}
$\vec{v} = \frac{\vec{\Delta x}}{\Delta t}$

En una dimensión:
$v = \frac{x_f - x_0}{t_f - t_0} = \frac{\Delta x}{\Delta t}$

\textbf{Velocidad instantánea:}
$v = \lim_{\Delta t \to 0} \frac{\Delta x}{\Delta t}$
\end{ecuacion}

\textbf{Diferencia:}
\begin{itemize}
    \item Rapidez usa distancia total
    \item Velocidad usa desplazamiento
    \item Rapidez siempre es positiva o cero
    \item Velocidad puede ser positiva, negativa o cero
\end{itemize}
\end{seccion}

\newpage

\section{ACELERACIÓN}
\begin{seccion}

\begin{definicion}
    \textbf{Aceleración:} Magnitud vectorial que mide el cambio de velocidad.
\end{definicion}

\begin{ecuacion}
\textbf{Ecuación de aceleración:}
$\vec{a} = \frac{\vec{\Delta v}}{\Delta t} = \frac{\vec{v_f} - \vec{v_0}}{t_f - t_0}$

En una dimensión:
$a = \frac{v_f - v_0}{\Delta t}$

\textbf{Aceleración instantánea:}
$a = \lim_{\Delta t \to 0} \frac{\Delta v}{\Delta t}$
\end{ecuacion}

\textbf{Unidades:} m/s², km/h², etc.

\vspace{0.3cm}

\textbf{Interpretación del signo:}
\begin{itemize}
    \item $a > 0$: el objeto acelera (aumenta velocidad en dirección positiva)
    \item $a < 0$: el objeto desacelera o frena
    \item $a = 0$: velocidad constante
\end{itemize}
\end{seccion}

\subsection{Esquema conceptual}

\begin{center}
\begin{tikzpicture}[scale=1.0]
% Título
\node[draw, fill=azuloscuro, text=white, rectangle, rounded corners, minimum width=6cm] at (0,3.5) {\textbf{ACELERACION}};

\node[font=\Large] at (0,2.5) {$a = \frac{\Delta v}{\Delta t}$};

\node[draw, fill=verdecorrecto!30, rectangle, minimum width=4cm, minimum height=0.8cm] at (-4,0.8) {$a > 0$ Acelera};
\node[font=\small, align=center] at (-4,0) {Velocidad aumenta};

\node[draw, fill=yellow!30, rectangle, minimum width=4cm, minimum height=0.8cm] at (0,0.8) {$a = 0$ MRU};
\node[font=\small, align=center] at (0,0) {Velocidad constante};

\node[draw, fill=red!30, rectangle, minimum width=4cm, minimum height=0.8cm] at (4,0.8) {$a < 0$ Frena};
\node[font=\small, align=center] at (4,0) {Velocidad disminuye};

% Flechas

\end{tikzpicture}
\end{center}

\newpage

\section{MOVIMIENTO RECTILÍNEO UNIFORME (MRU)}

\begin{seccion}
\begin{definicion}
    \textbf{Definición:} Movimiento en línea recta con velocidad constante.
\end{definicion}
\vspace{0.3cm}

\textbf{Características:}
\begin{itemize}
    \item Velocidad constante: $v = \text{constante}$
    \item Aceleración nula: $a = 0$
    \item Recorre distancias iguales en tiempos iguales
\end{itemize}

\begin{ecuacion}
\textbf{Ecuación general del movimiento:}
$x_f = x_0 + v_0 t + \frac{1}{2}at^2$

\textbf{Para MRU} ($a = 0$):
$x_f = x_0 + vt$

Si partimos del origen ($x_0 = 0$):
$x = vt$

Despejando para velocidad:
$v = \frac{x}{t}$
\end{ecuacion}
\end{seccion}


\subsection{Gráficas del MRU}

\begin{seccion}
\textbf{Gráfica posición vs tiempo:}

\begin{center}
\begin{tikzpicture}
\begin{axis}[
    width=12cm,
    height=6cm,
    axis lines=middle,
    xlabel={$t$ (s)},
    ylabel={$x$ (m)},
    xmin=0, xmax=6,
    ymin=0, ymax=25,
    xtick={0,1,2,3,4,5},
    ytick={0,5,10,15,20},
    grid=major,
    thick,
    every axis x label/.style={at={(current axis.right of origin)},anchor=west},
    every axis y label/.style={at={(current axis.above origin)},anchor=south}
]
% Línea recta con pendiente positiva (v > 0)
\addplot[color=azuloscuro, very thick, domain=0:5] {4*x};
\node[anchor=west] at (axis cs:3,15) {$v > 0$};
\end{axis}
\end{tikzpicture}
\end{center}

\begin{itemize}
    \item Línea recta
    \item Pendiente = velocidad
    \item Pendiente positiva: movimiento en dirección positiva
    \item Pendiente negativa: movimiento en dirección negativa
    \item Pendiente cero: reposo
\end{itemize}
\end{seccion}

\newpage

\begin{seccion}
\textbf{Gráfica velocidad vs tiempo:}

\begin{center}
\begin{tikzpicture}
\begin{axis}[
    width=12cm,
    height=6cm,
    axis lines=middle,
    xlabel={$t$ (s)},
    ylabel={$v$ (m/s)},
    xmin=0, xmax=6,
    ymin=0, ymax=10,
    xtick={0,1,2,3,4,5},
    ytick={0,2,4,6,8},
    grid=major,
    thick,
    every axis x label/.style={at={(current axis.right of origin)},anchor=west},
    every axis y label/.style={at={(current axis.above origin)},anchor=south}
]
% Línea horizontal (velocidad constante)
\addplot[color=azuloscuro, very thick] coordinates {(0,6) (5,6)};
\node[anchor=south] at (axis cs:2.5,6) {$v = \text{constante}$};

% Sombreado del área (distancia)
\addplot[color=azulmedio, fill=azulmedio, opacity=0.3] coordinates {(0,0) (5,0) (5,6) (0,6)} -- cycle;
\node[anchor=center] at (axis cs:2.5,3) {Área = distancia};
\end{axis}
\end{tikzpicture}
\end{center}

\begin{itemize}
    \item Línea horizontal
    \item Área bajo la curva = distancia recorrida
\end{itemize}
\end{seccion}

\newpage

\section{MOVIMIENTO UNIFORMEMENTE ACELERADO (MUA)}

\begin{seccion}
\begin{definicion}
    \textbf{Definición:} Movimiento en línea recta con aceleración constante.
\end{definicion}

\vspace{0.3cm}

\textbf{Características:}
\begin{itemize}
    \item Aceleración constante: $a = \text{constante}$
    \item Velocidad cambia uniformemente
    \item La velocidad aumenta (o disminuye) la misma cantidad en tiempos iguales
\end{itemize}
\end{seccion}

\subsection{Ecuaciones del MUA}

\begin{ecuacion}
Partimos de la ecuación general:
$$x_f = x_0 + v_0 t + \frac{1}{2}at^2$$

\vspace{0.3cm}

\textbf{Ecuación 1 - Posición en función del tiempo:}
$$x_f = x_0 + v_0 t + \frac{1}{2}at^2$$

\textbf{Ecuación 2 - Velocidad en función del tiempo:}
$$v_f = v_0 + at$$

\textbf{Ecuación 3 - Velocidad en función de la posición:}

Eliminando el tiempo:
$$v_f^2 = v_0^2 + 2a(x_f - x_0)$$

O simplificado:
$$v_f^2 = v_0^2 + 2a\Delta x$$

\textbf{Ecuación 4 - Distancia con velocidad media:}
$$x = \left(\frac{v_0 + v_f}{2}\right) t$$
\end{ecuacion}

\newpage

\subsection{Gráficas del MUA}

\begin{seccion}
\textbf{Gráfica posición vs tiempo:}

\begin{center}
\begin{tikzpicture}
\begin{axis}[
    width=12cm,
    height=7cm,
    axis lines=middle,
    xlabel={$t$ (s)},
    ylabel={$x$ (m)},
    xmin=0, xmax=5,
    ymin=0, ymax=50,
    xtick={0,1,2,3,4},
    ytick={0,10,20,30,40},
    grid=major,
    thick,
    every axis x label/.style={at={(current axis.right of origin)},anchor=west},
    every axis y label/.style={at={(current axis.above origin)},anchor=south}
]
% Parábola (a > 0)
\addplot[color=azuloscuro, very thick, domain=0:4.5] {2.5*x^2};
\node[anchor=west] at (axis cs:3,35) {$a > 0$};
\end{axis}
\end{tikzpicture}
\end{center}

\begin{itemize}
    \item Parábola
    \item Curvatura indica aceleración
    \item Cóncava hacia arriba si $a > 0$
    \item Cóncava hacia abajo si $a < 0$
\end{itemize}
\end{seccion}

\newpage

\begin{seccion}
\textbf{Gráfica velocidad vs tiempo:}

\begin{center}
\begin{tikzpicture}
\begin{axis}[
    width=12cm,
    height=7cm,
    axis lines=middle,
    xlabel={$t$ (s)},
    ylabel={$v$ (m/s)},
    xmin=0, xmax=6,
    ymin=0, ymax=25,
    xtick={0,1,2,3,4,5},
    ytick={0,5,10,15,20},
    grid=major,
    thick,
    every axis x label/.style={at={(current axis.right of origin)},anchor=west},
    every axis y label/.style={at={(current axis.above origin)},anchor=south}
]
% Línea inclinada (aceleración constante)
\addplot[color=azuloscuro, very thick, domain=0:5] {4*x};
\node[anchor=south west] at (axis cs:3,12) {pendiente = $a$};

% Sombreado del área (distancia)
\addplot[color=azulmedio, fill=azulmedio, opacity=0.3] coordinates {(0,0) (5,0) (5,20)} -- cycle;
\node[anchor=center] at (axis cs:2.5,6) {Área = distancia};
\end{axis}
\end{tikzpicture}
\end{center}

\begin{itemize}
    \item Línea recta inclinada
    \item Pendiente = aceleración
    \item Área bajo la curva = distancia recorrida
\end{itemize}
\end{seccion}

\newpage

\section{CAÍDA LIBRE}

\begin{seccion}

\begin{definicion}
    \textbf{Definición:} Movimiento vertical bajo la influencia de la gravedad.
\end{definicion}

\vspace{0.3cm}

\textbf{Características:}
\begin{itemize}
    \item Aceleración constante: $g = 9.8 m/s^2 ; g \approx 10 m/s^2$ (hacia abajo)
    \item Es un caso de MUA
    \item Dirección vertical
\end{itemize}

\begin{ecuacion}
\textbf{Ecuaciones de caída libre:}

Usamos las ecuaciones de MUA con $a = g$:

\vspace{0.3cm}

\textbf{Velocidad final:}
$v_f = v_0 + gt$

\textbf{Altura en función del tiempo:}
$y = y_0 + v_0 t + \frac{1}{2}gt^2$

Si se suelta desde el reposo ($v_0 = 0$) y desde $y_0 = h$:
$h = \frac{1}{2}gt^2$

\textbf{Velocidad en función de la altura:}
$v_f^2 = v_0^2 + 2gh$
\end{ecuacion}
\end{seccion}


\section{TIRO VERTICAL}

\begin{seccion}

\begin{definicion}
    \textbf{Definición:} Objeto lanzado verticalmente hacia arriba o hacia abajo.
\end{definicion}

\vspace{0.3cm}

\textbf{Características:}
\begin{itemize}
    \item Aceleración: $a = -g$ (siempre hacia abajo)
    \item En la subida: la velocidad disminuye
    \item En la altura máxima: $v = 0$
    \item En la bajada: la velocidad aumenta
\end{itemize}

\begin{ecuacion}
\textbf{Ecuaciones del tiro vertical:}

\textbf{Velocidad en cualquier instante:}
$v_f = v_0 - gt$

\textbf{Altura en función del tiempo:}
$y = y_0 + v_0 t - \frac{1}{2}gt^2$

\textbf{Altura máxima:}

En el punto más alto, $v_f = 0$:
$h_{max} = \frac{v_0^2}{2g}$

\textbf{Tiempo de subida:}
$t_{subida} = \frac{v_0}{g}$

\textbf{Tiempo de vuelo} (regresa al mismo nivel):
$t_{total} = \frac{2v_0}{g}$

\textbf{Velocidad en función de la altura:}
$v_f^2 = v_0^2 - 2gy$
\end{ecuacion}
\end{seccion}

\newpage

\section{TIRO PARABÓLICO}

\begin{seccion}

\begin{definicion}
    \textbf{Definición:} Movimiento en dos dimensiones bajo la influencia de la gravedad.
\end{definicion}

\vspace{0.3cm}

\textbf{Características:}
\begin{itemize}
    \item Combinación de MRU horizontal y MUA vertical
    \item Movimiento horizontal: velocidad constante
    \item Movimiento vertical: aceleración $g$ hacia abajo
    \item Trayectoria: parábola
\end{itemize}

\begin{ecuacion}
\textbf{Componentes de la velocidad inicial:}
$v_{0x} = v_0 \cos \theta$
$v_{0y} = v_0 \sin \theta$

\textbf{Ecuaciones horizontales (MRU):}
$x = v_{0x} t = v_0 \cos \theta \cdot t$
$v_x = v_{0x} = \text{constante}$

\textbf{Ecuaciones verticales (MUA con $a = -g$):}
$y = v_{0y} t - \frac{1}{2}gt^2 = v_0 \sin \theta \cdot t - \frac{1}{2}gt^2$
$v_y = v_{0y} - gt = v_0 \sin \theta - gt$

\textbf{Altura máxima:}
$h_{max} = \frac{v_0^2 \sin^2 \theta}{2g}$

\textbf{Tiempo de vuelo} (regresa al mismo nivel):
$t_{vuelo} = \frac{2v_0 \sin \theta}{g}$

\textbf{Alcance horizontal:}
$R = \frac{v_0^2 \sin 2\theta}{g}$

El alcance máximo ocurre cuando $\theta = 45°$
\end{ecuacion}
\end{seccion}

\subsection{Esquema conceptual}

\begin{center}
\begin{tikzpicture}[scale=0.85]
% Ejes
\draw[->, thick] (0,0) -- (8,0) node[right] {$x$};
\draw[->, thick] (0,0) -- (0,4) node[above] {$y$};

% Trayectoria parabólica
\draw[very thick, azuloscuro] (0,0) parabola bend (4,3) (8,0);

% Velocidad inicial
\draw[->, ultra thick, verdecorrecto] (0,0) -- (2,1.5) node[above right] {$\vec{v_0}$};
\draw[dashed] (0,0) -- (2,0) node[below, midway] {$v_{0x}$};
\draw[dashed] (2,0) -- (2,1.5) node[right, midway] {$v_{0y}$};
\draw (0.8,0) arc (0:37:0.8) node[right, midway] {$\theta$};

% Punto más alto
\filldraw[red] (4,3) circle (2pt) node[above] {$h_{max}$};
\draw[dashed] (4,0) -- (4,3);

% Alcance
\draw[<->, thick] (0,-0.5) -- (8,-0.5) node[midway, below] {Alcance: $R$};

% Componentes del movimiento
\node[draw, fill=fondoverde, align=center, minimum width=2.5cm] at (2,5) {Horizontal MRU};
\node[draw, fill=yellow!30, align=center, minimum width=2.5cm] at (6,5) {Vertical MUA};

% Gravedad
\draw[->, ultra thick, red] (7,3.5) -- (7,2) node[right, midway] {$g$};
\end{tikzpicture}
\end{center}


\newpage

\section{FORMULARIO - CINEMÁTICA}

\begin{tcolorbox}[colback=fondoverde, colframe=verdecorrecto, title=\textbf{Ecuación general del movimiento}]
$x_f = x_0 + v_0 t + \frac{1}{2}at^2$
\end{tcolorbox}

\begin{multicols}{2}

\subsection{Magnitudes fundamentales}

\begin{ecuacion}
\textbf{Velocidad:}
$v = \frac{\Delta x}{\Delta t}$

\textbf{Rapidez:}
$r = \frac{d}{t}$

\textbf{Aceleración:}
$a = \frac{\Delta v}{\Delta t}$
\end{ecuacion}

\subsection{Movimiento Rectilíneo Uniforme}

\begin{ecuacion}
$v = \frac{x}{t}$
$x = vt$
\end{ecuacion}

\subsection{Movimiento Uniformemente Acelerado}

\begin{ecuacion}
$v_f = v_0 + at$

$x = v_0 t + \frac{1}{2}at^2$

$v_f^2 = v_0^2 + 2ax$

$x = \left(\frac{v_0 + v_f}{2}\right) t$
\end{ecuacion}

\subsection{Caída Libre}

\begin{ecuacion}
$v_f = v_0 + gt$

$h = v_0 t + \frac{1}{2}gt^2$

$v_f^2 = v_0^2 + 2gh$

\textbf{Desde el reposo:}
$h = \frac{1}{2}gt^2$
\end{ecuacion}

\columnbreak

\subsection{Tiro Vertical}

\begin{ecuacion}
$v_f = v_0 - gt$

$y = v_0 t - \frac{1}{2}gt^2$

$v_f^2 = v_0^2 - 2gy$

\textbf{Altura máxima:}
$h_{max} = \frac{v_0^2}{2g}$

\textbf{Tiempo de subida:}
$t_{subida} = \frac{v_0}{g}$

\textbf{Tiempo total:}
$t_{total} = \frac{2v_0}{g}$
\end{ecuacion}

\subsection{Tiro Parabólico}

\begin{ecuacion}
\textbf{Componentes:}
$v_{0x} = v_0 \cos \theta$
$v_{0y} = v_0 \sin \theta$

\textbf{Horizontal:}
$x = v_0 \cos \theta \cdot t$

\textbf{Vertical:}
$y = v_0 \sin \theta \cdot t - \frac{1}{2}gt^2$

\textbf{Altura máxima:}
$h_{max} = \frac{v_0^2 \sin^2 \theta}{2g}$

\textbf{Tiempo de vuelo:}
$t_{vuelo} = \frac{2v_0 \sin \theta}{g}$

\textbf{Alcance:}
$R = \frac{v_0^2 \sin 2\theta}{g}$


\end{ecuacion}

\subsection{Constantes}

\begin{ecuacion}
$g \approx 10 m/s^2$
$g = 9.8 m/s^2$
\end{ecuacion}

\end{multicols}
